\documentclass[10pt,a4paper,twocolumn]{article}

% ── Pakete ────────────────────────────────────────────────────
\usepackage[T1]{fontenc}
\usepackage[utf8]{inputenc}
\usepackage[ngerman]{babel}
\usepackage{lmodern}
\usepackage[left=18mm,right=18mm,top=22mm,bottom=25mm]{geometry}
\usepackage{amsmath,amssymb,amsthm,mathtools}
\usepackage{booktabs}
\usepackage{float}
\usepackage{enumitem}
\usepackage[ruled,vlined,linesnumbered,german]{algorithm2e}
\usepackage{xcolor}
\usepackage{fancyhdr}
\usepackage{titlesec}
\usepackage{hyperref}
\usepackage{cleveref}

% ── Farben ────────────────────────────────────────────────────
\definecolor{darkblue}{RGB}{0,51,102}
\definecolor{noteblue}{RGB}{230,240,255}
\definecolor{warnred}{RGB}{255,235,235}

% ── Hyperref ──────────────────────────────────────────────────
\hypersetup{
    colorlinks=true,
    linkcolor=darkblue,
    citecolor=darkblue,
    urlcolor=darkblue,
}

% ── Seitenlayout ──────────────────────────────────────────────
\pagestyle{fancy}
\fancyhf{}
\fancyhead[L]{\small\textit{Toleranzintervalle -- Theorie und Implementierung}}
\fancyhead[R]{\small\thepage}
\renewcommand{\headrulewidth}{0.4pt}

% ── Abschnittsformatierung ────────────────────────────────────
\titleformat{\section}{\large\bfseries\color{darkblue}}{\thesection}{1em}{}
\titleformat{\subsection}{\normalsize\bfseries\color{darkblue}}{\thesubsection}{1em}{}

% ── Theorem-Umgebungen ────────────────────────────────────────
\theoremstyle{definition}
\newtheorem{definition}{Definition}[section]
\newtheorem{satz}[definition]{Satz}
\newtheorem{bemerkung}[definition]{Bemerkung}

% ── Hinweis-Boxen ─────────────────────────────────────────────
\newcommand{\hinweis}[1]{%
    \begin{center}
    \colorbox{noteblue}{\parbox{0.92\linewidth}{\small\textbf{Hinweis:} #1}}
    \end{center}
}
\newcommand{\warnung}[1]{%
    \begin{center}
    \colorbox{warnred}{\parbox{0.92\linewidth}{\small\textbf{Warnung:} #1}}
    \end{center}
}

% ── Abkürzungen ───────────────────────────────────────────────
\newcommand{\TI}{Toleranzintervall}
\newcommand{\KI}{Konfidenzintervall}
\newcommand{\R}{\mathbb{R}}
\newcommand{\E}{\mathbb{E}}
\newcommand{\Prob}{\mathbb{P}}
\newcommand{\Var}{\operatorname{Var}}
\newcommand{\xbar}{\bar{x}}

% ══════════════════════════════════════════════════════════════
\begin{document}
% ══════════════════════════════════════════════════════════════

\twocolumn[{%
\begin{center}
    {\LARGE\bfseries\color{darkblue}
    Toleranzintervalle:\\[3pt]
    Theorie, Verfahren und automatisierte Berechnung}\\[12pt]
    {\large Frank Zimmermann \quad\textbar\quad Technische Dokumentation}\\[6pt]
    {\normalsize Februar 2026}\\[4pt]
    \rule{0.6\textwidth}{0.5pt}
\end{center}
\vspace{6mm}

\begin{center}
\parbox{0.92\textwidth}{\small\textbf{Zusammenfassung.}
Dieses Dokument beschreibt die mathematischen Grundlagen und die
softwaretechnische Implementierung eines automatisierten
Toleranzintervall-Rechners. Ausgehend von der Abgrenzung zu
Konfidenzintervallen und Prädiktionsintervallen werden vier Verfahren
entwickelt: verteilungsfrei (Ordnungsstatistiken), normalverteilt
($k$-Faktor), lognormalverteilt (Transformation) und
Weibull-basiert (parametrischer Bootstrap). Für das verteilungsfreie
Verfahren wird eine Optimierung vorgestellt, die durch Verschieben
eines gleitenden Fensters über die Ordnungsstatistiken das engste
Intervall bei gegebener Konfidenz findet -- eine Verbesserung
gegenüber dem in der Standardliteratur (Meeker et~al., 2017)
beschriebenen symmetrischen Ansatz. Ein Entscheidungsalgorithmus
wählt automatisch das geeignete Verfahren anhand von Stichprobenumfang und
Verteilungstests aus. Die Implementierung erfolgt in Python mit
PySide6-Oberfläche.}
\end{center}
\vspace{8mm}
}]

\tableofcontents
\vspace{6mm}

% ══════════════════════════════════════════════════════════════
\section{Warum Toleranzintervalle?}
\label{sec:motivation}
% ══════════════════════════════════════════════════════════════

In der industriellen Praxis stehen Ingenieure und Qualitätssicherer
regelmäßig vor der Frage: \emph{Welchen Bereich decken die Messwerte
meiner Produktion ab?} Diese Frage klingt einfach, aber die Antwort
hängt davon ab, was genau man wissen möchte.

\subsection{Drei Arten von Intervallen}

Gegeben seien $n$ unabhängige Messwerte $x_1, \ldots, x_n$ aus einer
Grundgesamtheit mit unbekannter Verteilung $F$.

\begin{definition}[\KI{} für den Mittelwert]
Ein $100(1-\alpha)\%$-\KI{} ist ein Intervall $[L, U]$, sodass
\begin{equation}
    \Prob(\mu \in [L, U]) \geq 1 - \alpha
\end{equation}
wobei $\mu = \E[X]$ der wahre Populationsmittelwert ist.
Das \KI{} macht eine Aussage über einen \emph{Parameter}.
\end{definition}

\begin{definition}[Prädiktionsintervall]
Ein $100(1-\alpha)\%$-Prädiktionsintervall ist ein Intervall $[L, U]$,
sodass ein \emph{einzelner zukünftiger} Messwert $X_{n+1}$ mit
Wahrscheinlichkeit $\geq 1-\alpha$ in $[L, U]$ fällt.
\end{definition}

\begin{definition}[\TI{}]
\label{def:ti}
Ein $100p / 100(1-\alpha)$-\TI{} ist ein Intervall $[L, U]$,
das \emph{mindestens den Anteil $p$ der Grundgesamtheit} überdeckt,
und zwar mit einer Konfidenz von $1 - \alpha$. Formal:
\begin{equation}
    \Prob\!\Big(\Prob(X \in [L, U]) \geq p\Big) \geq 1 - \alpha
    \label{eq:ti_def}
\end{equation}
\end{definition}

Der entscheidende Unterschied: Das \KI{} schätzt einen Parameter
(Punkt), das Prädiktionsintervall sagt einen einzelnen Wert voraus,
das \TI{} charakterisiert einen \emph{Anteil der Population}.

\subsection{Praktische Motivation}

Toleranzintervalle beantworten Fragen wie:

\begin{itemize}[nosep,leftmargin=*]
    \item \emph{Welcher Bereich umfasst 95\% aller produzierten Teile?}
    \item \emph{Können wir mit 99\% Konfidenz sagen, dass 90\% der
          Lebensdauern über 1000\,h liegen?}
    \item \emph{Liegen unsere Messwerte innerhalb der Spezifikation?}
\end{itemize}

In der industriellen Qualitätssicherung nach ISO~16269 sind
Toleranzintervalle das Standardwerkzeug, um aus Stichproben auf
die Streuung der Grundgesamtheit zu schließen.

\hinweis{Ein 95/95-\TI{} bedeutet: Mit 95\% Konfidenz überdeckt
das Intervall mindestens 95\% der Population. Beides -- Abdeckung
\emph{und} Konfidenz -- muss angegeben werden.}


% ══════════════════════════════════════════════════════════════
\section{Verteilungsfreies Toleranzintervall}
\label{sec:verteilungsfrei}
% ══════════════════════════════════════════════════════════════

\subsection{Grundidee}

Das verteilungsfreie Verfahren macht \emph{keine Annahme} über die
Form der Grundgesamtheit. Es nutzt ausschließlich die
Ordnungsstatistiken der Stichprobe.

\begin{definition}[Ordnungsstatistiken]
Seien $x_1, \ldots, x_n$ die Messwerte. Die Ordnungsstatistiken sind
die aufsteigend sortierten Werte:
\[
    x_{(1)} \leq x_{(2)} \leq \cdots \leq x_{(n)}
\]
\end{definition}

\subsection{Zweiseitiges Intervall}

Für das Intervall $[x_{(\ell)}, x_{(u)}]$ mit $1 \leq \ell < u \leq n$
ist die Wahrscheinlichkeit, dass dieses Intervall mindestens den
Anteil $p$ der Population überdeckt, gegeben durch die
\emph{Coverage Probability of a Tolerance Interval}
(Meeker et~al., 2017, Formel~5.9):

\begin{satz}[CPTI-Formel]
\label{satz:cpti}
Die Konfidenz für das Intervall $[x_{(\ell)}, x_{(u)}]$ ist
\[
    \mathrm{CPTI}(n, \ell, u, p) = \sum_{i=0}^{d-1} \binom{n}{i}\,p^i\,(1-p)^{n-i}
\]
wobei $d = u - \ell$ der \emph{Span} (Anzahl der Werte zwischen den Grenzen)
ist. Dies entspricht der Binomial-CDF: $\mathrm{CPTI} = F_{\mathrm{Binom}}(d-1;\,n,\,p)$.
\end{satz}

\hinweis{Die Konfidenz hängt \textbf{nur vom Span $d = u - \ell$ ab},
nicht von der Position $(\ell, u)$ innerhalb der sortierten Stichprobe.
Diese Eigenschaft ist der Schlüssel zur weiter unten beschriebenen Optimierung
(\cref{sec:optimierung}).}

Für den Spezialfall $\ell=1, u=n$ (äußerste Ordnungsstatistiken,
$d = n-1$) vereinfacht sich die Formel:

\begin{satz}[Konfidenz für Minimum und Maximum]
\label{satz:minmax}
Für den Spezialfall $d = n-1$ (Intervall $[x_{(1)}, x_{(n)}]$) gilt:
\[
    C(n, p) = 1 - n\,p^{n-1} + (n-1)\,p^n
\]
\end{satz}

\begin{proof}
Folgt direkt aus \cref{satz:cpti} mit $d = n-1$: Die Summation
über die Binomialterme ergibt nach Umformung die geschlossene Form
$1 - np^{n-1} + (n-1)p^n$.
\end{proof}

Das \emph{minimale $n$} für vorgegebene Abdeckung $p$ und
Konfidenz $1-\alpha$ ergibt sich durch Lösen von
\begin{equation}
    1 - n\,p^{n-1} + (n-1)\,p^n \geq 1 - \alpha
    \label{eq:min_n_two}
\end{equation}
nach dem kleinsten ganzzahligen $n$.

\subsection{Symmetrisches Trimmen nach Meeker}
\label{sec:meeker_symm}

Der klassische Ansatz (Meeker et~al., 2017, Kap.~5) bestimmt einen
Parameter $\nu$ aus Tabelle~J.11, der angibt, wie viele
Ordnungsstatistiken von den Rändern entfernt werden. Die Aufteilung
erfolgt symmetrisch:
\[
    \ell = \lfloor \nu/2 \rfloor, \qquad u = n - \lceil \nu/2 \rceil + 1
\]
Für den Fall, dass die resultierende Konfidenz den geforderten Wert
$1-\alpha$ deutlich übersteigt, beschreibt Meeker eine
\emph{Interpolationsmethode} (Example~5.10): Man bildet einen
gewichteten Durchschnitt zwischen dem konservativen und dem
nichtkonservativen Randpaar.

Dieser Ansatz hat zwei Nachteile:
\begin{enumerate}[leftmargin=*]
    \item \textbf{Symmetrische Aufteilung:} Bei schiefen Verteilungen
    liegt die Datenspannweite nicht symmetrisch um das Zentrum. Das
    symmetrische Trimmen kann ein unnötig breites Intervall ergeben.
    \item \textbf{Interpolation:} Die Interpolation zwischen konservativem
    und nichtkonservativem Intervall ist nicht mehr streng verteilungsfrei
    und kann das Intervall sogar \emph{breiter} machen als das konservative
    Ergebnis (siehe \cref{sec:vergleich_meeker}).
\end{enumerate}

\subsection{Optimierung: Gleitendes Fenster}
\label{sec:optimierung}

Da die Konfidenz gemäß \cref{satz:cpti} ausschließlich vom Span $d$
abhängt, kann man das Fenster $[\ell, \ell + d]$ über alle zulässigen
Positionen $\ell = 1, \ldots, n - d$ verschieben und dasjenige wählen,
für das die Intervallbreite $x_{(\ell+d)} - x_{(\ell)}$ minimal wird.

\begin{satz}[Optimales verteilungsfreies TI]
\label{satz:opt_df}
Sei $d^*$ der minimale Span mit
$F_{\mathrm{Binom}}(d^* - 1;\, n,\, p) \geq 1 - \alpha$,
und sei
\[
    \ell^* = \arg\min_{\ell \in \{1,\ldots,n-d^*\}}
        \big(x_{(\ell + d^*)} - x_{(\ell)}\big).
\]
Dann ist $[x_{(\ell^*)},\; x_{(\ell^* + d^*)}]$ das engste
verteilungsfreie $(p,\,1{-}\alpha)$-Toleranzintervall mit der
Konfidenzgarantie:
\[
    \mathrm{CPTI}(n, \ell^*, \ell^* + d^*, p) = F_{\mathrm{Binom}}(d^* - 1;\, n,\, p) \geq 1-\alpha
\]
\end{satz}

\begin{proof}
Die Konfidenz ist invariant unter Verschiebung des Fensters (da sie nur
von $d$ abhängt). Also hat das verschobene Intervall exakt dieselbe
statistische Garantie wie jedes andere Intervall mit gleichem $d$.
Die Wahl von $\ell^*$ minimiert lediglich die Breite -- ein reiner
Informationsgewinn ohne Verlust der Konfidenz.
\end{proof}

Der Algorithmus läuft in $O(n)$ -- ein einziger Durchlauf über die
sortierten Daten genügt. Die Implementierung ist trivial:

\begin{algorithm}
\caption{Optimales verteilungsfreies TI}
\label{alg:opt_df}
\SetKwInOut{Input}{Eingabe}
\SetKwInOut{Output}{Ausgabe}
\Input{Sortierte Daten $x_{(1)}, \ldots, x_{(n)}$; Abdeckung $p$; Konfidenz $1-\alpha$}
\Output{Optimales Intervall $[x_{(\ell^*)}, x_{(u^*)}]$}
Finde $d^*$ als kleinstes $d$ mit $F_{\text{Binom}}(d-1; n, p) \geq 1-\alpha$\;
$w^* \leftarrow \infty$; $\ell^* \leftarrow 1$\;
\For{$\ell \leftarrow 1$ \KwTo $n - d^*$}{
    $w \leftarrow x_{(\ell + d^*)} - x_{(\ell)}$\;
    \lIf{$w < w^*$}{$w^* \leftarrow w$; $\ell^* \leftarrow \ell$}
}
\Return{$[x_{(\ell^*)}, x_{(\ell^* + d^*)}]$}
\end{algorithm}

\subsection{Vergleich mit Meeker: Compound-Amount-Daten}
\label{sec:vergleich_meeker}

Als Testdaten verwenden wir die 100~Messwerte aus Meeker et~al.\
(2017, Table~5.1): geordnete Konzentrationen (ppm) eines chemischen
Prozesses, mit $x_{(1)} = 1.49$ und $x_{(100)} = 58.11$.
Die Daten sind deutlich rechtsschief.

Gesucht: zweiseitiges TI mit $p = 0.90$, $1-\alpha = 0.95$.

\begin{table}[H]
\centering
\caption{Vergleich der Verfahren für die Meeker-Daten ($n=100$, $p=0.90$, $1-\alpha = 0.95$)}
\label{tab:meeker_vergleich}
\small
\begin{tabular}{@{}lcccc@{}}
\toprule
Methode & $\ell$ & $u$ & Konfidenz & Breite \\
\midrule
Naiv ($x_{(1)}, x_{(n)}$) & 1 & 100 & 99.97\% & 56.62 \\
Meeker Ex.\,5.9 (symm.) & 2 & 98 & 97.63\% & 35.66 \\
Meeker Ex.\,5.10 (interp.) & -- & -- & $\approx$\,95\% & 37.78 \\
\textbf{Optimiert (unser)} & \textbf{1} & \textbf{97} & \textbf{97.63\%} & \textbf{31.75} \\
\bottomrule
\end{tabular}
\end{table}

Die Ergebnisse zeigen:

\begin{enumerate}[leftmargin=*]
    \item Der \textbf{naive Ansatz} ($\ell=1, u=100$) liefert zwar eine
    Konfidenz von 99.97\%, ist aber mit einer Breite von 56.62 stark
    überkonservativ -- die statistisch nicht genutzte Reservekonfidenz
    ist verschwendete Information.

    \item \textbf{Meeker Ex.\,5.9} trimmt symmetrisch ($\nu = 5 \Rightarrow \ell=2, u=98$)
    und reduziert den Span auf $d = 96$ (CPTI $= 0.9763$). Das Intervall
    $[1.66,\; 37.32]$ ist deutlich enger.

    \item \textbf{Meeker Ex.\,5.10} interpoliert zwischen dem konservativen
    und nichtkonservativen Intervall, um die Konfidenz auf genau 95\% zu
    treffen. Paradoxerweise wird das Intervall dadurch \emph{breiter}
    (37.78 statt 35.66), weil die Interpolation die obere Grenze
    Richtung $x_{(99)} = 53.43$ zieht.

    \item \textbf{Unsere Optimierung} nutzt denselben Span $d = 96$ wie
    Meeker und hat identische Konfidenz (97.63\%), verschiebt aber das
    Fenster nach $\ell = 1, u = 97$. Die obere Grenze fällt von 37.32
    auf 33.24 -- eine \textbf{Reduktion der Intervallbreite um 11\%}.
\end{enumerate}

\hinweis{Bei symmetrischen Daten liefert die Optimierung dasselbe
Ergebnis wie Meekers symmetrischer Ansatz. Der Vorteil zeigt sich
gerade bei schiefen Verteilungen, wo die Extremwerte auf einer Seite
konzentriert sind -- also genau dann, wenn man das verteilungsfreie
Verfahren braucht, weil die Normalverteilung nicht passt.}

Warum verwendet Meeker nicht diese Optimierung? Der Grund ist
historisch: Das Tabellensystem (Table~J.11) wurde für die
Taschenrechner-Ära entworfen, in der man einen Parameter $\nu$
nachschlägt und symmetrisch aufteilt. Die Fenster-Optimierung
braucht Zugriff auf die gesamten sortierten Daten und eine
Schleife -- trivial für einen Rechner, aber nicht für Papier und
Bleistift. In der zweiten Auflage (2017) wurde das Framework
beibehalten, obwohl Meeker R verwendet.

\subsection{Einseitiges Intervall}

Für ein oberes einseitiges \TI{} $(-\infty, x_{(n)}]$ gilt:
\begin{equation}
    \Prob\!\big(F(x_{(n)}) \geq p\big) = 1 - p^n
    \label{eq:df_one}
\end{equation}

Das minimale $n$ folgt aus $1 - p^n \geq 1 - \alpha$:
\begin{equation}
    n \geq \frac{\ln \alpha}{\ln p}
    \label{eq:min_n_one}
\end{equation}

Analog für das untere einseitige Intervall $[x_{(1)}, \infty)$.

\subsection{Mindeststichprobengrößen}

\Cref{tab:min_n} zeigt die resultierenden Mindest-$n$ für gängige
Kombinationen von $p$ und $1-\alpha$.

\begin{table}[H]
\centering
\caption{Minimales $n$ für verteilungsfreie TI
($r=1, s=n$)}
\label{tab:min_n}
\small
\begin{tabular}{@{}cccc@{}}
\toprule
$p$ & $1-\alpha$ & einseitig & zweiseitig \\
\midrule
0.90 & 0.90 &  22 &  38 \\
0.90 & 0.95 &  29 &  46 \\
0.90 & 0.99 &  44 &  64 \\
\midrule
0.95 & 0.90 &  45 &  77 \\
0.95 & 0.95 &  59 &  93 \\
0.95 & 0.99 &  90 & 130 \\
\midrule
0.99 & 0.90 & 230 & 388 \\
0.99 & 0.95 & 299 & 473 \\
0.99 & 0.99 & 459 & 662 \\
\bottomrule
\end{tabular}
\end{table}

\warnung{Für $p=0.99$ bei $1-\alpha=0.99$ sind \textbf{662 Messwerte}
nötig. Bei teuren Prüfungen (z.\,B.\ Lebensdauertests) ist das
verteilungsfreie Verfahren oft nicht praktikabel -- man muss dann
auf ein parametrisches Verfahren ausweichen.}


% ══════════════════════════════════════════════════════════════
\section{Normalverteilungs-Toleranzintervall}
\label{sec:normal}
% ══════════════════════════════════════════════════════════════

\subsection{Voraussetzung}

Die Messwerte stammen aus einer Normalverteilung:
$X_i \sim \mathcal{N}(\mu, \sigma^2)$ mit unbekanntem $\mu$ und
$\sigma$.

\subsection{Der $k$-Faktor}

Das \TI{} hat die Form:
\begin{equation}
    [\xbar - k \cdot s, \quad \xbar + k \cdot s]
    \label{eq:normal_ti}
\end{equation}
wobei $\xbar$ der Stichprobenmittelwert und
$s = \sqrt{\frac{1}{n-1}\sum_{i=1}^n (x_i - \xbar)^2}$
die Stichprobenstandardabweichung ist.

Der Faktor $k$ muss so bestimmt werden, dass \cref{eq:ti_def}
erfüllt ist.

\subsubsection{Einseitiger Fall}

Für ein oberes einseitiges \TI{} $(-\infty, \xbar + k \cdot s]$
folgt $k$ aus der \emph{nichtzentralen $t$-Verteilung}:

\begin{satz}[$k$-Faktor, einseitig]
\label{satz:k_one}
Der einseitige $k$-Faktor erfüllt:
\begin{equation}
    k = \frac{t_{n-1,\,1-\alpha,\,\delta}}{\sqrt{n}}
    \label{eq:k_one}
\end{equation}
wobei $t_{n-1,\,1-\alpha,\,\delta}$ das $(1-\alpha)$-Quantil der
nichtzentralen $t$-Verteilung mit $n-1$ Freiheitsgraden und
Nichtzentralitätsparameter $\delta = z_p \sqrt{n}$ ist, und
$z_p = \Phi^{-1}(p)$ das $p$-Quantil der Standardnormalverteilung.
\end{satz}

Die Herleitung nutzt, dass
$T = \frac{\xbar - \mu}{s/\sqrt{n}}$ einer $t$-Verteilung mit
$n-1$ Freiheitsgraden folgt und die Toleranzforderung auf ein
Quantilproblem der nichtzentralen $t$-Verteilung zurückgeführt
werden kann.

\subsubsection{Zweiseitiger Fall}

Der zweiseitige Fall ist analytisch komplexer. Eine gängige
Approximation nach \textsc{Howe} (1969) lautet:

\begin{satz}[$k$-Faktor, zweiseitig -- Howe-Approximation]
\label{satz:k_two}
\begin{equation}
    k = z_{(1+p)/2} \cdot
    \sqrt{\frac{(n-1)\,(1 + 1/n)}{\chi^2_{1-\alpha,\,n-1}}}
    \label{eq:k_two}
\end{equation}
wobei $\chi^2_{1-\alpha,\,n-1}$ das $(1-\alpha)$-Quantil der
$\chi^2$-Verteilung mit $n-1$ Freiheitsgraden ist (unteres Quantil).
\end{satz}

Die Idee: Die Unsicherheit in $\xbar$ wird durch den Faktor
$(1 + 1/n)$ berücksichtigt, die Unsicherheit in $s^2$ durch die
$\chi^2$-Verteilung.

\subsection{Verhalten des $k$-Faktors}

Für $n \to \infty$ konvergiert $k$ gegen $z_{(1+p)/2}$, d.\,h.\
die Unsicherheit in der Schätzung von $\mu$ und $\sigma$ verschwindet.
Für kleines $n$ ist $k$ deutlich größer, was die
Schätzunsicherheit kompensiert. \Cref{tab:k_values} zeigt einige
Werte.

\begin{table}[H]
\centering
\caption{$k$-Faktoren für $p=0.95$, $1-\alpha=0.95$ (zweiseitig,
Howe-Approximation)}
\label{tab:k_values}
\small
\begin{tabular}{@{}cc@{}}
\toprule
$n$ & $k$ \\
\midrule
5   & 4.4141 \\
10  & 3.0793 \\
15  & 2.7523 \\
20  & 2.5836 \\
30  & 2.4183 \\
50  & 2.2689 \\
100 & 2.1330 \\
$\infty$ & 1.9600 \\
\bottomrule
\end{tabular}
\end{table}


% ══════════════════════════════════════════════════════════════
\section{Normalitätstest: Shapiro-Wilk}
\label{sec:shapiro}
% ══════════════════════════════════════════════════════════════

Bevor das Normal-\TI{} angewendet wird, muss die Normalverteilungs-Annahme
geprüft werden. Der Shapiro-Wilk-Test ist für kleine bis mittlere
Stichproben ($n \leq 5000$) einer der mächtigsten Tests.

\subsection{Teststatistik}

\begin{definition}[Shapiro-Wilk-Teststatistik]
\begin{equation}
    W = \frac{\left(\sum_{i=1}^{n} a_i\, x_{(i)}\right)^2}
             {\sum_{i=1}^{n} (x_i - \xbar)^2}
    \label{eq:shapiro}
\end{equation}
wobei $x_{(i)}$ die Ordnungsstatistiken sind und die Gewichte
$a_i$ aus der Kovarianzmatrix der Ordnungsstatistiken der
Standardnormalverteilung berechnet werden:
\begin{equation}
    (a_1, \ldots, a_n) = \frac{\mathbf{m}^T \mathbf{V}^{-1}}
    {\|\mathbf{V}^{-1}\mathbf{m}\|}
\end{equation}
mit $\mathbf{m} = \E[\mathbf{Z}_{(i)}]$ den erwarteten
Ordnungsstatistiken und $\mathbf{V}$ deren Kovarianzmatrix.
\end{definition}

\subsection{Testentscheidung}

\begin{itemize}[nosep,leftmargin=*]
    \item $H_0$: Die Daten stammen aus einer Normalverteilung
    \item $H_1$: Die Daten stammen nicht aus einer Normalverteilung
    \item Ablehnung von $H_0$, wenn der $p$-Wert $< \alpha$
          (typisch: $\alpha = 0.05$)
\end{itemize}

\hinweis{Der Shapiro-Wilk-Test testet auf \emph{Abweichung} von der
Normalität. Ein hoher $p$-Wert \emph{bestätigt} nicht die
Normalverteilung -- er sagt nur, dass die Daten nicht signifikant
davon abweichen. Bei kleinem $n$ hat der Test wenig Macht.}


% ══════════════════════════════════════════════════════════════
\section{Lognormal-Toleranzintervall}
\label{sec:lognormal}
% ══════════════════════════════════════════════════════════════

\subsection{Voraussetzung}

Die Messwerte $x_i > 0$ folgen einer Lognormalverteilung, d.\,h.\
$\ln(x_i) \sim \mathcal{N}(\mu_{\ln}, \sigma_{\ln}^2)$.

Typische Anwendungen: Konzentrationen, Einkommen, Partikelgrößen,
bestimmte Lebensdauern.

\subsection{Verfahren}

Die Grundidee ist denkbar einfach: Man transformiert die Daten
auf die Log-Skala, berechnet dort ein Normal-\TI{} und
transformiert zurück.

\begin{enumerate}[nosep,leftmargin=*]
    \item Berechne $y_i = \ln(x_i)$
    \item Berechne $\bar{y}$ und $s_y$ auf der Log-Skala
    \item Bestimme $k$ nach \cref{eq:k_two} (oder \cref{eq:k_one})
    \item Das \TI{} auf der Originalskala ist:
\end{enumerate}

\begin{equation}
    \big[\exp(\bar{y} - k \cdot s_y), \quad
          \exp(\bar{y} + k \cdot s_y)\big]
    \label{eq:lognormal_ti}
\end{equation}

\subsection{Verteilungstest}

Die Lognormalitäts-Annahme wird geprüft, indem der Shapiro-Wilk-Test
auf die \emph{log-transformierten} Daten $y_i = \ln(x_i)$ angewendet
wird.

\hinweis{Das Lognormal-\TI{} ist \emph{asymmetrisch} auf der
Originalskala -- die obere Grenze kann sehr viel weiter vom Median
entfernt sein als die untere. Das ist physikalisch sinnvoll für
rechtsschiefe Daten.}


% ══════════════════════════════════════════════════════════════
\section{Weibull-Toleranzintervall}
\label{sec:weibull}
% ══════════════════════════════════════════════════════════════

\subsection{Die Weibull-Verteilung}

Die Weibull-Verteilung ist \emph{das} Standardmodell für Lebensdauer-
und Zuverlässigkeitsanalysen. Ihre Dichte lautet:

\begin{equation}
    f(x; c, \lambda) = \frac{c}{\lambda}
    \left(\frac{x}{\lambda}\right)^{c-1}
    \exp\!\left[-\left(\frac{x}{\lambda}\right)^c\right]
    \label{eq:weibull_pdf}
\end{equation}
für $x > 0$, mit Formparameter $c > 0$ und Skalenparameter $\lambda > 0$.

Die Verteilungsfunktion ist:
\begin{equation}
    F(x) = 1 - \exp\!\left[-\left(\frac{x}{\lambda}\right)^c\right]
    \label{eq:weibull_cdf}
\end{equation}

\subsection{Interpretation des Formparameters}

Der Parameter $c$ bestimmt die Form der Ausfallrate:
\begin{itemize}[nosep,leftmargin=*]
    \item $c < 1$: abnehmende Ausfallrate (Frühausfälle, \emph{infant mortality})
    \item $c = 1$: konstante Ausfallrate (Exponentialverteilung)
    \item $c > 1$: zunehmende Ausfallrate (Alterung, Verschleiß)
    \item $c \approx 3.6$: annähernd normalverteilt
\end{itemize}

\subsection{Parameterschätzung}

Die Parameter $c$ und $\lambda$ werden per
\emph{Maximum-Likelihood-Estimation} (MLE) geschätzt.
Die Log-Likelihood ist:

\begin{equation}
    \ell(c, \lambda) = n \ln c - n c \ln \lambda
    + (c-1)\sum_{i=1}^n \ln x_i
    - \sum_{i=1}^n \left(\frac{x_i}{\lambda}\right)^c
    \label{eq:weibull_loglik}
\end{equation}

Die Maximierung erfolgt numerisch (z.\,B.\ über
\texttt{scipy.stats.weibull\_min.fit}).

\subsection{Anpassungstest}

Da \texttt{scipy.stats.anderson} die Weibull-Verteilung nicht direkt
unterstützt, verwenden wir einen \emph{Kolmogorow-Smirnow-Test mit
Monte-Carlo-Korrektur} (Lilliefors-Ansatz):

\begin{enumerate}[nosep,leftmargin=*]
    \item Fitte Weibull$(c, \lambda)$ per MLE an die Daten
    \item Berechne die KS-Statistik:
        \begin{equation}
            D_n = \sup_x \big|F_n(x) - F(x; \hat{c}, \hat{\lambda})\big|
            \label{eq:ks_stat}
        \end{equation}
    \item Simuliere $N_{\text{MC}} = 500$ Stichproben aus
          Weibull$(\hat{c}, \hat{\lambda})$
    \item Für jede Simulation: fitte neu und berechne $D_n^{(j)}$
    \item Der Monte-Carlo-$p$-Wert ist:
        \begin{equation}
            p_{\text{MC}} = \frac{1}{N_{\text{MC}}}
            \sum_{j=1}^{N_{\text{MC}}} \mathbf{1}[D_n^{(j)} \geq D_n]
            \label{eq:mc_pvalue}
        \end{equation}
\end{enumerate}

\warnung{Der naive KS-$p$-Wert (ohne Monte-Carlo-Korrektur) ist
\emph{zu liberal}, wenn die Parameter aus denselben Daten geschätzt
werden. Die Korrektur via Simulation ist essenziell
(Lilliefors-Problem).}

\subsection{Toleranzintervall via parametrischem Bootstrap}

Für die Weibull-Verteilung existiert keine geschlossene Formel für
den $k$-Faktor analog zum Normalfall. Stattdessen verwenden wir
\emph{parametrischen Bootstrap}:

\begin{enumerate}[nosep,leftmargin=*]
    \item Fitte Weibull$(c, \lambda)$ an die Originaldaten
    \item Für $b = 1, \ldots, B$ (typisch $B = 2000$):
    \begin{enumerate}[nosep]
        \item Generiere $n$ Werte aus Weibull$(\hat{c}, \hat{\lambda})$
        \item Fitte neu: $\hat{c}^{(b)}, \hat{\lambda}^{(b)}$
        \item Berechne die Zielquantile:
            \begin{align}
                q_L^{(b)} &= F^{-1}\!\left(\frac{1-p}{2};\,
                    \hat{c}^{(b)}, \hat{\lambda}^{(b)}\right) \\
                q_U^{(b)} &= F^{-1}\!\left(\frac{1+p}{2};\,
                    \hat{c}^{(b)}, \hat{\lambda}^{(b)}\right)
            \end{align}
    \end{enumerate}
    \item Die Konfidenzgrenzen des \TI{}s sind:
        \begin{align}
            L &= Q_{\alpha/2}\!\left(\{q_L^{(b)}\}_{b=1}^B\right) \\
            U &= Q_{1-\alpha/2}\!\left(\{q_U^{(b)}\}_{b=1}^B\right)
            \label{eq:bootstrap_bounds}
        \end{align}
        wobei $Q_\gamma$ das $\gamma$-Quantil der Bootstrap-Verteilung bezeichnet.
\end{enumerate}

Die Idee: Die Bootstrap-Verteilung der Quantilschätzer bildet die
Parameterunsicherheit ab. Indem wir konservative Perzentile wählen
(unteres Ende der unteren Grenze, oberes Ende der oberen), erhalten
wir ein \TI{} mit der geforderten Konfidenz.

\subsection{Weibull-Quantilfunktion}

Die Quantilfunktion der Weibull-Verteilung hat eine geschlossene Form:
\begin{equation}
    F^{-1}(q) = \lambda \cdot \big(-\ln(1-q)\big)^{1/c}
    \label{eq:weibull_quantile}
\end{equation}

Dies macht die Bootstrap-Berechnung numerisch effizient.


% ══════════════════════════════════════════════════════════════
\section{Automatische Methodenwahl}
\label{sec:algorithmus}
% ══════════════════════════════════════════════════════════════

Der zentrale Beitrag der Implementierung ist ein Entscheidungs-Algorithmus, der automatisch das geeignetste Verfahren auswählt.

\subsection{Entscheidungslogik}

\begin{algorithm}
\SetAlgoLined
\KwIn{Daten $x_1, \ldots, x_n$; Abdeckung $p$; Konfidenz $1-\alpha$;
      Seite; Shapiro-$\alpha_s$}
\KwOut{Toleranzintervall $[L, U]$ mit Methode}

\tcc{Schritt 1--3: Parametrische Verfahren (nutzen alle $n$ Datenpunkte)}
\If{$n \geq 3$}{
    $(W, p_W) \leftarrow$ Shapiro-Wilk$(x_1, \ldots, x_n)$\;
    \If{$p_W \geq \alpha_s$}{
        \Return Normal-TI ($k$-Faktor)\;
    }
}
\If{$\forall i: x_i > 0$ \textbf{und} $n \geq 3$}{
    $(W, p_W) \leftarrow$ Shapiro-Wilk$(\ln x_1, \ldots, \ln x_n)$\;
    \If{$p_W \geq \alpha_s$}{
        \Return Lognormal-TI ($k$-Faktor auf Log-Skala)\;
    }
}
\If{$\forall i: x_i > 0$ \textbf{und} $n \geq 5$}{
    $(c, \lambda, D_n, p_{\text{MC}}) \leftarrow$ Weibull-GoF\;
    \If{$p_{\text{MC}} \geq \alpha_s$}{
        \Return Weibull-TI (parametrischer Bootstrap)\;
    }
}
\tcc{Schritt 4: Verteilungsfrei (braucht viele Daten)}
$n_{\min} \leftarrow$ \texttt{min\_n}$(p, 1-\alpha, \text{Seite})$\;
\If{$n \geq n_{\min}$}{
    \Return Verteilungsfreies TI (Ordnungsstatistiken, Fensteroptimierung)\;
}
\tcc{Schritt 5: Letzter Ausweg}
\Return Normal-TI mit Warnung (Fallback)\;

\caption{Automatische Methodenwahl (parametrisch zuerst)}
\label{algo:auto}
\end{algorithm}

\subsection{Begründung der Reihenfolge}

Die Priorisierung folgt dem Prinzip \emph{maximaler
Informationsnutzung}:

\begin{enumerate}[nosep,leftmargin=*]
    \item \textbf{Parametrische Verfahren zuerst} (Normal, Lognormal,
          Weibull): Diese nutzen \emph{alle $n$ Datenpunkte}
          ($\xbar$, $s$, bzw.\ MLE-Schätzer) und liefern daher
          engere Intervalle. Besonders bei kleinem $n$ ist der
          Informationsgewinn gegenüber dem verteilungsfreien Verfahren
          erheblich.
    \item \textbf{Normal vor Lognormal}: Vermeidet, dass bei
          annähernd symmetrischen positiven Daten fälschlicherweise
          Lognormal gewählt wird, nur weil Shapiro-Wilk auf
          $\ln(x)$ zufällig einen höheren $p$-Wert liefert.
    \item \textbf{Weibull nach Lognormal}: Die $k$-Faktor-Methode
          ist analytisch sauberer und schneller als der
          Monte-Carlo-Bootstrap.
    \item \textbf{Verteilungsfrei als bewusste Wahl}: Kommt erst
          zum Zug, wenn keine parametrische Verteilung passt
          \emph{und} $n \geq n_{\min}$. Bei $n \gg n_{\min}$
          kann die Fensteroptimierung (\cref{sec:optimierung}) innere
          Ordnungsstatistiken wählen und liefert dann ein
          konkurrenzfähiges Intervall.
    \item \textbf{Fallback}: Wenn weder eine Verteilung passt noch
          genug Daten für das verteilungsfreie Verfahren vorliegen,
          liefert das Normal-TI zumindest eine Orientierung --
          aber mit deutlicher Warnung.
\end{enumerate}

\warnung{Das verteilungsfreie Verfahren ist \textbf{nicht} der
konservativste oder sicherste Ansatz -- bei $n \approx n_{\min}$
ist es der \emph{schwächste}, weil es nur die Extremwerte
$x_{(1)}$ und $x_{(n)}$ verwenden kann. Ein parametrisches
Verfahren, dessen Verteilungsannahme durch den Shapiro-Wilk-Test
gestützt wird, liefert engere und informativere Intervalle.}


% ══════════════════════════════════════════════════════════════
\section{Einseitige Toleranzintervalle}
\label{sec:einseitig}
% ══════════════════════════════════════════════════════════════

Alle beschriebenen Verfahren unterstützen einseitige Intervalle:

\subsection{Oberes Toleranzintervall}
$(-\infty, U]$: Garantiert, dass mindestens $p$ der Population
\emph{unterhalb} von $U$ liegt.
\begin{itemize}[nosep,leftmargin=*]
    \item Verteilungsfrei: $U = x_{(n)}$, Konfidenz $1 - p^n$
    \item Normal: $U = \xbar + k_{\text{eins}} \cdot s$
          mit $k$ nach \cref{eq:k_one}
\end{itemize}

\subsection{Unteres Toleranzintervall}
$[L, \infty)$: Garantiert, dass mindestens $p$ der Population
\emph{oberhalb} von $L$ liegt.
\begin{itemize}[nosep,leftmargin=*]
    \item Verteilungsfrei: $L = x_{(1)}$
    \item Normal: $L = \xbar - k_{\text{eins}} \cdot s$
\end{itemize}

\subsection{Anwendung}

Einseitige Intervalle sind relevant, wenn nur eine Richtung
interessiert:
\begin{itemize}[nosep,leftmargin=*]
    \item Mindest-Lebensdauer: \emph{unteres} TI
    \item Maximal-Schadstoffgehalt: \emph{oberes} TI
    \item Mindest-Festigkeit: \emph{unteres} TI
\end{itemize}

Der einseitige $k$-Faktor ist kleiner als der zweiseitige
(bei gleichem $p$ und $1-\alpha$), weil die
,,Wahrscheinlichkeitsmasse'' nicht auf zwei Seiten verteilt werden
muss.


% ══════════════════════════════════════════════════════════════
\section{Verifikation und Tests}
\label{sec:tests}
% ══════════════════════════════════════════════════════════════

Die Implementierung wird durch 21~Unit-Tests abgesichert, die
alle wesentlichen Eigenschaften prüfen.

\subsection{Tests der Mindeststichprobengrößen}

Die berechneten Werte werden gegen die Referenztabellen geprüft:
\begin{align*}
    n_{\min}(p{=}0.95, 1{-}\alpha{=}0.95, \text{zweiseitig}) &= 93 \\
    n_{\min}(p{=}0.95, 1{-}\alpha{=}0.95, \text{einseitig}) &= 59 \\
    n_{\min}(p{=}0.99, 1{-}\alpha{=}0.99, \text{zweiseitig}) &= 662 \\
    n_{\min}(p{=}0.90, 1{-}\alpha{=}0.90, \text{zweiseitig}) &= 38
\end{align*}

\subsection{Tests des $k$-Faktors}

\begin{itemize}[nosep,leftmargin=*]
    \item \textbf{Positivität}: $k > 0$ für alle gültigen Parameter
    \item \textbf{Monotonie in $n$}: $k(n=20) > k(n=100)$ --
          mehr Daten $\Rightarrow$ kleinerer $k$-Faktor
    \item \textbf{Einseitig $<$ Zweiseitig}: $k_{\text{eins}} < k_{\text{zwei}}$
          bei gleichen Parametern
\end{itemize}

\subsection{Tests der Methodenwahl}

Die automatische Entscheidungslogik wird mit synthetischen Daten
geprüft:

\begin{table}[H]
\centering
\caption{Testszenarien für die automatische Methodenwahl}
\label{tab:tests}
\small
\begin{tabular}{@{}lll@{}}
\toprule
Szenario & $n$ & Erwartete Methode \\
\midrule
$\mathcal{N}(\mu, \sigma^2)$ & 100 & Normal ($k$-Faktor) \\
$\mathcal{N}(\mu, \sigma^2)$ &  20 & Normal ($k$-Faktor) \\
$\text{Lognormal}(\mu, \sigma)$ &  25 & Lognormal \\
Bimodal & 100 & Verteilungsfrei \\
Bimodal &  15 & Fallback (Warnung) \\
Mit NaN-Werten & 100 & Korrekte $n$-Zählung \\
$n < 2$ & 1 & ValueError \\
\bottomrule
\end{tabular}
\end{table}

\subsection{Tests der Intervallgrenzen}

\begin{itemize}[nosep,leftmargin=*]
    \item \textbf{Verteilungsfrei}: Grenzen = $x_{(1)}$ und $x_{(n)}$;
          tatsächliche Konfidenz $\geq$ geforderte Konfidenz
    \item \textbf{Normal}: Intervall symmetrisch um $\xbar$, d.\,h.\
          $\xbar - L = U - \xbar$ (bis Maschinengenauigkeit)
    \item \textbf{Lognormal}: Beide Grenzen strikt positiv; $L < U$;
          negative Eingabedaten $\Rightarrow$ \texttt{ValueError}
    \item \textbf{Weibull}: Grenzen strikt positiv;
          negative Eingabedaten $\Rightarrow$ \texttt{ValueError}
    \item \textbf{Einseitig}: Korrekte \texttt{None}-Werte für die
          offene Seite
\end{itemize}


% ══════════════════════════════════════════════════════════════
\section{Praktische Hinweise}
\label{sec:praxis}
% ══════════════════════════════════════════════════════════════

\subsection{Stichprobengröße}

Die Faustregel: \emph{Mehr Daten $=$ engere Intervalle}.
Aber die Verbesserung ist nichtlinear -- der Sprung von $n=10$
auf $n=30$ bringt viel mehr als von $n=100$ auf $n=120$.

Für das verteilungsfreie Verfahren ist $n$ eine harte Schwelle:
Entweder man hat genug Messwerte oder nicht.
Für die parametrischen Verfahren gibt es keine Mindestschwelle,
aber bei $n < 10$ sind die Intervalle so breit, dass sie kaum
noch nützlich sind.

\subsection{Umgang mit Ausreißern}

Toleranzintervalle sind \emph{nicht robust} gegen Ausreißer:
\begin{itemize}[nosep,leftmargin=*]
    \item Verteilungsfrei: Wird direkt von Extremwerten bestimmt
    \item Normal: $\xbar$ und $s$ sind beide ausreißerempfindlich
    \item Lognormal: Etwas robuster durch die Log-Transformation
    \item Weibull: MLE kann durch Ausreißer verzerrt werden
\end{itemize}

\warnung{Ausreißer sollten \emph{vor} der TI-Berechnung
untersucht werden. Einen echten Messfehler zu entfernen ist
gerechtfertigt; einen unbequemen, aber validen Messwert zu
streichen ist Betrug.}

\subsection{Wahl der Abdeckung und Konfidenz}

Typische Kombinationen in der Praxis:
\begin{itemize}[nosep,leftmargin=*]
    \item \textbf{90/90}: Screening, Voruntersuchungen
    \item \textbf{95/95}: Standard in der Qualitätssicherung (ISO~16269)
    \item \textbf{99/95}: Sicherheitskritische Anwendungen
    \item \textbf{99/99}: Luftfahrt, Medizinprodukte
\end{itemize}

\subsection{Grenzen des Fallback-Verfahrens}

Wenn keine der getesteten Verteilungen passt, liefert das
Normal-TI als Fallback ein Ergebnis. Dieses ist jedoch mit
Vorsicht zu genießen:
\begin{itemize}[nosep,leftmargin=*]
    \item Bei schiefen Verteilungen kann die untere Grenze
          negativ werden (z.\,B.\ bei Lebensdauerdaten)
    \item Die tatsächliche Abdeckung kann deutlich von $p$
          abweichen
    \item Die Empfehlung lautet: mehr Daten sammeln und/oder
          domänenspezifisches Wissen über die Verteilung einbringen
\end{itemize}


% ══════════════════════════════════════════════════════════════
\section{Mathematische Zusammenfassung}
\label{sec:zusammenfassung}
% ══════════════════════════════════════════════════════════════

\Cref{tab:formeln} fasst die zentralen Formeln aller vier
Verfahren zusammen.

\begin{table*}[t]
\centering
\caption{Übersicht der Verfahren und Formeln für zweiseitige Toleranzintervalle}
\label{tab:formeln}
\small
\begin{tabular}{@{}p{2.8cm}p{3.5cm}p{4.5cm}p{3.5cm}@{}}
\toprule
\textbf{Verfahren} & \textbf{Intervall} & \textbf{Schlüsselformel} & \textbf{Voraussetzung} \\
\midrule
Normal ($k$-Faktor)
    & $[\xbar - ks,\; \xbar + ks]$
    & $k = z_{(1+p)/2}\sqrt{\frac{(n-1)(1+1/n)}{\chi^2_{1-\alpha,n-1}}}$
    & Shapiro-Wilk $p \geq \alpha_s$ \\[4pt]
Lognormal
    & $[e^{\bar{y}-ks_y},\; e^{\bar{y}+ks_y}]$
    & $k$ wie Normal, auf $y_i = \ln x_i$
    & $x_i > 0$, SW auf $\ln x$ \\[4pt]
Weibull (Bootstrap)
    & $[Q_{\alpha/2}(q_L^{(b)}),$ $Q_{1-\alpha/2}(q_U^{(b)})]$
    & $q^{(b)} = \hat\lambda^{(b)}(-\ln(1-q))^{1/\hat c^{(b)}}$
    & $x_i > 0$, MC-KS $p \geq \alpha_s$ \\[4pt]
Verteilungsfrei
    & $[x_{(\ell^*)},\; x_{(\ell^*+d^*)}]$
    & CPTI $= F_{\text{Binom}}(d^*{-}1;\,n,\,p)$
    & $n \geq n_{\min}$, kein param.\ Fit \\
\bottomrule
\end{tabular}
\end{table*}


% ══════════════════════════════════════════════════════════════
\section{Implementierung}
\label{sec:implementierung}
% ══════════════════════════════════════════════════════════════

\subsection{Architektur}

Die Software besteht aus zwei Schichten:

\begin{description}[nosep,leftmargin=*,font=\bfseries]
    \item[Rechenkern] (\texttt{tolerance\_intervals.py}):
        Reine Berechnungslogik ohne GUI-Abhängigkeiten.
        Kann auch standalone als Python-Modul verwendet werden.
    \item[GUI] (\texttt{ti\_gui.py}):
        PySide6-Oberfläche mit Dateneingabe, Parameterkonfiguration
        und Ergebnisanzeige.
\end{description}

\subsection{Abhängigkeiten}

\begin{itemize}[nosep,leftmargin=*]
    \item \textbf{NumPy}: Datenverarbeitung, Ordnungsstatistiken
    \item \textbf{SciPy}: Verteilungen, $k$-Faktor (nichtzentrale
          $t$-Verteilung, $\chi^2$), Shapiro-Wilk-Test,
          Weibull-MLE
    \item \textbf{PySide6}: GUI-Framework (Qt für Python)
    \item \textbf{PyInstaller}: Standalone-Binary-Erzeugung
\end{itemize}

\subsection{Projektstruktur}

\begin{verbatim}
toleranzintervall-rechner/
  main.py              # Einstiegspunkt
  build.py             # Build-Skript
  src/
    tolerance_intervals.py
    ti_gui.py
  tests/
    test_tolerance.py
  .run/                # PyCharm Configs
  pyproject.toml
  requirements.txt
  requirements-dev.txt
\end{verbatim}


% ══════════════════════════════════════════════════════════════
\section{Durchgerechnete Beispiele}
\label{sec:beispiele}
% ══════════════════════════════════════════════════════════════

Zum besseren Verständnis werden hier alle vier Verfahren an
konkreten Daten vorgeführt.

\subsection{Beispiel 1: Normal -- Messwerte einer Fertigung}

Gegeben seien $n = 15$ Durchmesser-Messwerte (in mm):
\[
    4.2,\; 4.5,\; 4.1,\; 4.8,\; 4.3,\; 4.6,\; 4.4,\; 4.7,\;
    4.2,\; 4.5,\; 4.3,\; 4.6,\; 4.4,\; 4.5,\; 4.3
\]

\textbf{Schritt 1: Verteilungsfrei möglich?}
Für $p = 0.95$, $1-\alpha = 0.95$, zweiseitig gilt
$n_{\min} = 93$. Da $n = 15 < 93$: nein.

\textbf{Schritt 2: Normalverteilung?}
Shapiro-Wilk: $W = 0.9726$, $p = 0.8951 \geq 0.05$
$\Rightarrow$ Normalverteilung wird nicht abgelehnt.

\textbf{Schritt 3: $k$-Faktor berechnen.}
Mit $n = 15$, $p = 0.95$, $1-\alpha = 0.95$ (zweiseitig):
\begin{align*}
    z_{0.975} &= 1.9600 \\[2pt]
    \chi^2_{0.05,\,14} &= 6.5706 \\[2pt]
    k &= 1.9600 \cdot \sqrt{\frac{14 \cdot (1 + 1/15)}{6.5706}} = 2.9548
\end{align*}

\textbf{Schritt 4: Intervall berechnen.}
\begin{align*}
    \xbar &= 4.4267 \qquad s = 0.1981 \\[2pt]
    L &= 4.4267 - 2.9548 \times 0.1981 = 3.8414 \\
    U &= 4.4267 + 2.9548 \times 0.1981 = 5.0120
\end{align*}

\textbf{Ergebnis:} $[3.841,\; 5.012]$ mm.
Mit 95\% Konfidenz liegen mindestens 95\% aller Durchmesser
in diesem Intervall.

\subsection{Beispiel 2: Verteilungsfrei -- Meeker-Referenzdaten}

Gegeben seien die $n = 100$ geordneten Konzentrationsmesswerte (ppm)
aus Meeker et~al.\ (2017, Table~5.1) mit
$x_{(1)} = 1.49$, $x_{(97)} = 33.24$, $x_{(98)} = 37.32$,
$x_{(100)} = 58.11$.

Gesucht: $p = 0.90$, $1-\alpha = 0.95$, zweiseitig.

\textbf{Schritt 1:} Minimaler Span. Finde kleinstes $d$ mit
$F_{\text{Binom}}(d-1;\,100,\,0.9) \geq 0.95$.
Lösung: $d^* = 96$ mit CPTI $= 0.9763$.

\textbf{Schritt 2:} Gleitendes Fenster. Verschiebe $[\ell, \ell + 96]$
über $\ell = 1, \ldots, 4$:
\begin{align*}
    \ell=1: &\quad [1.49,\; 33.24] \quad\text{Breite } 31.75 \quad\checkmark \\
    \ell=2: &\quad [1.66,\; 37.32] \quad\text{Breite } 35.66 \quad\text{(Meeker)} \\
    \ell=3: &\quad [2.05,\; 53.43] \quad\text{Breite } 51.38 \\
    \ell=4: &\quad [2.24,\; 58.11] \quad\text{Breite } 55.87
\end{align*}

Alle vier haben identische Konfidenz $97.63\%$, da $d = 96$ konstant ist.
Die Fenster-Optimierung wählt $\ell^* = 1$:

\textbf{Ergebnis:} $[1.49,\; 33.24]$ mit Konfidenz 97.63\%.
Meeker gibt $[1.66,\; 37.32]$ an -- 11\% breiter bei identischer Konfidenz.
Der Vorteil entsteht durch die Rechtsschiefe der Daten: die oberen
Extremwerte (53.43, 58.11) erzwingen beim symmetrischen Trimmen ein
unnötig breites Intervall.

\subsection{Beispiel 3: Lognormal -- Konzentrationsdaten}

Gegeben seien $n = 25$ Schadstoffkonzentrationen (mg/l),
die den Shapiro-Wilk-Test auf Normalität \emph{nicht} bestehen
($p = 0.0035$), aber auf der Log-Skala normalverteilt sind
($p = 0.4382$).

Auf der Log-Skala:
\begin{align*}
    \bar{y} &= 2.9622 \qquad s_y = 0.5234 \\[2pt]
    k &= z_{0.975} \cdot \sqrt{\frac{24 \cdot (1 + 1/25)}
         {\chi^2_{0.05,\,24}}} = 2.6313
\end{align*}

Rücktransformation:
\begin{align*}
    L &= \exp(2.9622 - 2.6313 \times 0.5234) = 4.879 \\
    U &= \exp(2.9622 + 2.6313 \times 0.5234) = 76.66
\end{align*}

\textbf{Ergebnis:} $[4.88,\; 76.66]$ mg/l.
Beachte die starke Asymmetrie: die obere Grenze liegt
$\approx 15\times$ weiter vom Median entfernt als die untere.
Das ist physikalisch sinnvoll für rechtsschiefe Konzentrationen.

\subsection{Beispiel 4: Weibull -- Lebensdauerdaten}

Gegeben seien $n = 30$ Lebensdauern (Stunden) eines Bauteils
mit abnehmender Ausfallrate. Shapiro-Wilk lehnt sowohl
Normalverteilung ($p = 0.0005$) als auch Lognormalverteilung
($p = 0.0045$) ab.

Weibull-Fit ergibt: $\hat{c} = 0.934$, $\hat{\lambda} = 675.4$~h.
Der Monte-Carlo-KS-Test: $p_{\text{MC}} = 0.492 \geq 0.05$
$\Rightarrow$ Weibull wird nicht abgelehnt.

Parametrischer Bootstrap mit $B = 2000$ liefert:
\[
    [L, U] = [3.47,\; 4153]~\text{h}
\]

\textbf{Zum Vergleich:} Das Normal-Fallback hätte
$[-70,\; 1876]$~h ergeben -- eine \emph{negative untere
Grenze} für Lebensdauern, was physikalisch sinnlos ist.
Das Weibull-TI liefert das korrekte Ergebnis: alle Grenzen
sind positiv.


% ══════════════════════════════════════════════════════════════
\section{Herleitung der $k$-Faktor-Formel}
\label{sec:herleitung}
% ══════════════════════════════════════════════════════════════

\subsection{Einseitiger Fall -- exakte Lösung}

Für ein oberes \TI{} $(-\infty, \xbar + ks]$ fordern wir:
\begin{equation}
    \Prob\!\left(\Phi\!\left(\frac{\xbar + ks - \mu}{\sigma}\right)
    \geq p\right) = 1 - \alpha
    \label{eq:k_derive_1}
\end{equation}

Sei $Z_p = \Phi^{-1}(p)$. Die Forderung
$\frac{\xbar + ks - \mu}{\sigma} \geq Z_p$ ist äquivalent zu:
\begin{equation}
    k \geq \frac{Z_p \sigma - (\xbar - \mu)}{s}
    = \frac{Z_p}{\frac{s}{\sigma}} - \frac{\xbar - \mu}{s}
\end{equation}

Definiere $T = \frac{\xbar - \mu}{s/\sqrt{n}} \sim t_{n-1}$
und verwende $\frac{(n-1)s^2}{\sigma^2} \sim \chi^2_{n-1}$.
Dann lässt sich zeigen, dass die Bedingung äquivalent ist zu:
\begin{equation}
    k\sqrt{n} \geq T^* \sim t_{n-1}(\delta)
\end{equation}
wobei $T^*$ einer nichtzentralen $t$-Verteilung mit
$\delta = Z_p\sqrt{n}$ folgt. Also:
\begin{equation}
    k = \frac{t_{n-1,\,1-\alpha,\,Z_p\sqrt{n}}}{\sqrt{n}}
\end{equation}

Dies ist eine \emph{exakte} Lösung, die numerisch über die
Quantilfunktion der nichtzentralen $t$-Verteilung berechnet wird.

\subsection{Zweiseitiger Fall -- Approximation}

Für das zweiseitige \TI{} $[\xbar - ks, \xbar + ks]$
fordern wir:
\begin{equation}
    \Prob\!\left(\Phi\!\left(\frac{\xbar + ks - \mu}{\sigma}\right)
    - \Phi\!\left(\frac{\xbar - ks - \mu}{\sigma}\right)
    \geq p\right) = 1 - \alpha
\end{equation}

Dies führt auf ein zweidimensionales Problem in $\xbar$ und $s$,
das keine geschlossene Lösung hat. Howes Approximation vereinfacht
durch die Faktorisierung:
\begin{itemize}[nosep,leftmargin=*]
    \item Die Unsicherheit in $\xbar$ wird durch $(1 + 1/n)$
          berücksichtigt (Varianz von $\xbar$)
    \item Die Unsicherheit in $s^2$ wird durch das $\chi^2$-Quantil
          eingefangen (Verteilung von $(n-1)s^2/\sigma^2$)
\end{itemize}

Die resultierende Formel
\begin{equation}
    k = z_{(1+p)/2} \cdot
    \sqrt{\frac{(n-1)(1 + 1/n)}{\chi^2_{1-\alpha,\,n-1}}}
\end{equation}
ist für $n \geq 10$ eine gute Approximation. Für kleinere $n$
gibt es tabellarische Korrekturen (siehe Meeker et al., 2017).


% ══════════════════════════════════════════════════════════════
\section{Vergleich der Verfahren}
\label{sec:vergleich}
% ══════════════════════════════════════════════════════════════

\subsection{Intervallbreite}

Bei gleichen Daten und gleichen Parametern ($p$, $1-\alpha$)
liefern die Verfahren unterschiedlich breite Intervalle.
Generell gilt:

\begin{center}
\begin{tabular}{@{}ll@{}}
\toprule
Eigenschaft & Rangfolge \\
\midrule
Engste Intervalle & Parametrisch (richtige Verteilung) \\
Breiteste Intervalle & Verteilungsfrei \\
Robustheit & Verteilungsfrei $>$ Parametrisch \\
Rechenzeit & Normal $\ll$ Weibull \\
\bottomrule
\end{tabular}
\end{center}

Das verteilungsfreie Intervall ist immer das konservativste
(breiteste), weil es keine Verteilungsannahme ausnutzt.
Im Gegenzug ist es immun gegen Modellfehler.

\subsection{Wann welches Verfahren?}

\begin{table}[H]
\centering
\caption{Empfehlungen für die Praxis}
\label{tab:empfehlung}
\small
\begin{tabular}{@{}p{3.2cm}p{4cm}@{}}
\toprule
\textbf{Situation} & \textbf{Empfehlung} \\
\midrule
$n > n_{\min}$, keine Verteilungsinfo &
    Verteilungsfrei \\[3pt]
$n < n_{\min}$, Daten symmetrisch &
    Normal \\[3pt]
$n < n_{\min}$, Daten rechtsschief, positiv &
    Lognormal oder Weibull \\[3pt]
Lebensdauerdaten, $c < 1.5$ &
    Weibull \\[3pt]
Keine Verteilung passt &
    Mehr Daten sammeln \\
\bottomrule
\end{tabular}
\end{table}

\subsection{Rechenaufwand}

\begin{itemize}[nosep,leftmargin=*]
    \item \textbf{Verteilungsfrei}: $O(n \log n)$ (Sortieren)
    \item \textbf{Normal}: $O(n)$ plus Quantilberechnung
    \item \textbf{Lognormal}: $O(n)$ wie Normal
    \item \textbf{Weibull GoF}: $O(N_{\text{MC}} \cdot n)$ für
          Monte-Carlo-KS
    \item \textbf{Weibull Bootstrap}: $O(B \cdot n)$ für $B$
          Bootstrap-Replikationen
\end{itemize}

In der Praxis dauert der Weibull-Pfad einige Sekunden
(ca.\ 2--5\,s für $n=30$, $B=2000$), alle anderen Verfahren
sind quasi instantan ($<$\,100\,ms).


% ══════════════════════════════════════════════════════════════
\section{Abgrenzung: TI vs.\ KI vs.\ Prädiktionsintervall}
\label{sec:abgrenzung}
% ══════════════════════════════════════════════════════════════

Ein häufiger Fehler in der Praxis ist die Verwechslung der drei
Intervalltypen. Am Beispiel von $n = 20$ normalverteilten
Messwerten mit $\xbar = 100$ und $s = 10$:

\begin{table}[H]
\centering
\caption{Drei Intervalltypen im Vergleich ($\alpha = 0.05$)}
\label{tab:vergleich}
\small
\begin{tabular}{@{}lcc@{}}
\toprule
\textbf{Intervalltyp} & \textbf{Formel} & \textbf{Breite} \\
\midrule
95\%-KI für $\mu$ &
    $\xbar \pm t_{0.975,19}\frac{s}{\sqrt{n}}$ &
    $[95.3,\; 104.7]$ \\[3pt]
95\%-Prädiktions-I. &
    $\xbar \pm t_{0.975,19}\,s\sqrt{1+\frac{1}{n}}$ &
    $[78.4,\; 121.6]$ \\[3pt]
95/95-TI &
    $\xbar \pm k \cdot s$ &
    $[72.5,\; 127.5]$ \\
\bottomrule
\end{tabular}
\end{table}

Das \KI{} ist mit Abstand das engste, weil es nur den
\emph{Mittelwert} einschließen muss -- einen einzelnen Punkt.
Das Prädiktionsintervall ist breiter, weil es einen
\emph{zukünftigen Einzelwert} abdecken muss. Das \TI{} ist
am breitesten, weil es \emph{95\% aller zukünftigen Werte}
gleichzeitig abdecken muss, und zwar mit einer zusätzlichen
Konfidenzsicherheit.

Die Breiten-Reihenfolge gilt immer:
\[
    |KI| \;<\; |PI| \;<\; |TI|
\]

\hinweis{Wer ein KI berechnet, aber ein TI braucht, unterschätzt
die Streuung der Population drastisch. Das ist einer der
häufigsten statistischen Fehler in technischen Berichten.}


% ══════════════════════════════════════════════════════════════
\section{Zusammenhang mit Prozessfähigkeitskennzahlen}
\label{sec:cpk}
% ══════════════════════════════════════════════════════════════

In der Qualitätssicherung werden häufig Prozessfähigkeitskennzahlen
wie $C_p$ und $C_{pk}$ verwendet. Es besteht ein enger Zusammenhang
zu Toleranzintervallen.

\subsection{Definitionen}

Gegeben Spezifikationsgrenzen $\text{LSL}$ (Lower Specification Limit)
und $\text{USL}$ (Upper Specification Limit):

\begin{align}
    C_p &= \frac{\text{USL} - \text{LSL}}{6\sigma}
    \label{eq:cp} \\[4pt]
    C_{pk} &= \min\!\left(
        \frac{\text{USL} - \mu}{3\sigma},\;
        \frac{\mu - \text{LSL}}{3\sigma}
    \right)
    \label{eq:cpk}
\end{align}

\subsection{Verbindung zu Toleranzintervallen}

Ein 99.73/50-\TI{} (Abdeckung 99.73\%, Konfidenz 50\%)
entspricht genau dem Bereich $\mu \pm 3\sigma$ -- dem
Nenner von $C_p$. Ein Prozess mit $C_p = 1$ hat also ein
natürliches Streuband, das gerade die Spezifikationsgrenzen
ausfüllt.

In der Praxis ist jedoch die \emph{Unsicherheit} in der Schätzung
von $\mu$ und $\sigma$ relevant. Hier bieten Toleranzintervalle
den entscheidenden Vorteil: Sie berücksichtigen die
Stichprobenunsicherheit explizit über die Konfidenz $1-\alpha$.

\begin{satz}[TI und Prozessfähigkeit]
Ein Prozess erfüllt die Spezifikation mit Abdeckung $p$ und
Konfidenz $1-\alpha$, wenn das entsprechende \TI{} vollständig
innerhalb $[\text{LSL}, \text{USL}]$ liegt:
\begin{equation}
    \xbar - k \cdot s \geq \text{LSL}
    \quad\text{und}\quad
    \xbar + k \cdot s \leq \text{USL}
    \label{eq:ti_spec}
\end{equation}
\end{satz}

Dies ist eine \emph{strengere} Forderung als $C_{pk} \geq 1$,
weil der $k$-Faktor größer ist als~3 (er berücksichtigt die
endliche Stichprobe). Erst für $n \to \infty$ konvergieren
beide Aussagen.

\hinweis{$C_{pk} \geq 1$ aus $n = 30$ Messwerten ist
\emph{keine} Garantie, dass der Prozess dauerhaft fähig ist.
Ein 95/95-TI innerhalb der Spezifikation ist die belastbarere
Aussage.}


% ══════════════════════════════════════════════════════════════
\section{Bootstrap-Verfahren im Detail}
\label{sec:bootstrap}
% ══════════════════════════════════════════════════════════════

\subsection{Warum Bootstrap?}

Für die Normalverteilung existiert eine analytische Lösung
($k$-Faktor). Für andere Verteilungen ist dies oft nicht der
Fall, insbesondere wenn:
\begin{itemize}[nosep,leftmargin=*]
    \item Die Pivotgröße $\frac{\hat{\theta} - \theta}{\hat{\text{se}}}$
          keine bekannte Verteilung hat
    \item Die Quantile der Verteilung keine einfache Funktion der
          Parameter sind
    \item Konfidenzgrenzen für Funktionale (nicht nur Parameter)
          gesucht sind
\end{itemize}

\subsection{Parametrischer vs.\ nichtparametrischer Bootstrap}

\begin{description}[nosep,leftmargin=*,font=\bfseries]
    \item[Parametrisch:] Simuliere aus der geschätzten Verteilung
        $\hat{F}(\cdot\,; \hat{\boldsymbol{\theta}})$.
        Nutzt die Verteilungsannahme aus, liefert engere Intervalle.
    \item[Nichtparametrisch:] Ziehe mit Zurücklegen aus der
        empirischen Verteilung $F_n$.
        Keine Verteilungsannahme, aber breitere Intervalle
        und instabil bei kleinem $n$.
\end{description}

Im Weibull-Fall verwenden wir den \emph{parametrischen} Bootstrap,
weil:
\begin{enumerate}[nosep,leftmargin=*]
    \item Die Weibull-Annahme bereits durch den GoF-Test gestützt ist
    \item Bei $n = 20$--$30$ der nichtparametrische Bootstrap
          für Extremquantile ($p = 0.025, 0.975$) unzuverlässig ist
    \item Die parametrische Variante die MLE-Eigenschaften
          (Konsistenz, Effizienz) nutzt
\end{enumerate}

\subsection{Konvergenzverhalten}

Die Genauigkeit des Bootstrap-\TI{}s hängt von der Anzahl
der Replikationen $B$ ab. Für die Perzentil-Methode gilt
approximativ:
\begin{equation}
    \text{se}(\hat{q}_\gamma) \approx
    \sqrt{\frac{\gamma(1-\gamma)}{B}} \cdot
    \frac{1}{f(\hat{q}_\gamma)}
    \label{eq:bootstrap_se}
\end{equation}
wobei $f$ die Dichte der Bootstrap-Verteilung ist.

Für $\gamma = 0.025$ und $B = 2000$ ergibt sich
$\text{se} \approx \frac{0.156}{f(\hat{q}_{0.025})} \cdot
\frac{1}{\sqrt{2000}} \approx \frac{0.0035}{f(\hat{q})}$,
was für praktische Zwecke ausreichend genau ist.

\warnung{Bei $B < 500$ können die Bootstrap-Konfidenzgrenzen
instabil sein. Bei sicherheitskritischen Anwendungen empfiehlt
sich $B \geq 5000$.}


% ══════════════════════════════════════════════════════════════
\section{Sensitivitätsanalyse}
\label{sec:sensitivitaet}
% ══════════════════════════════════════════════════════════════

\subsection{Einfluss des Signifikanzniveaus $\alpha_s$}

Die Wahl von $\alpha_s$ für den Shapiro-Wilk-Test beeinflusst
die Methodenwahl:

\begin{itemize}[nosep,leftmargin=*]
    \item \textbf{$\alpha_s = 0.01$ (streng)}: Normalverteilung
        wird seltener abgelehnt $\Rightarrow$ häufiger Normal-TI.
        Risiko: TI basiert auf falscher Verteilungsannahme.
    \item \textbf{$\alpha_s = 0.10$ (liberal)}: Normalverteilung
        wird häufiger abgelehnt $\Rightarrow$ häufiger
        Lognormal/Weibull/Fallback. Risiko: Unnötig komplexes
        Verfahren bei fast-normalen Daten.
\end{itemize}

Der Standardwert $\alpha_s = 0.05$ ist ein praxiserprobter
Kompromiss. Bei kleinem $n$ ($< 15$) hat der Shapiro-Wilk-Test
allerdings wenig Macht -- er lehnt fast nie ab, unabhängig von
der wahren Verteilung.

\subsection{Einfluss von $n$ auf die Intervallbreite}

Für das Normal-TI (zweiseitig, $p=0.95$, $1-\alpha=0.95$)
zeigt \Cref{tab:breite} die relative Intervallbreite
$2k \cdot s$ normiert auf $\sigma$:

\begin{table}[H]
\centering
\caption{Relative Breite des 95/95-Normal-TI}
\label{tab:breite}
\small
\begin{tabular}{@{}ccc@{}}
\toprule
$n$ & $k$ & $2k$ (in Einheiten $\sigma$) \\
\midrule
5   & 4.414 & 8.83 \\
10  & 3.079 & 6.16 \\
20  & 2.584 & 5.17 \\
50  & 2.269 & 4.54 \\
100 & 2.133 & 4.27 \\
$\infty$ & 1.960 & 3.92 \\
\bottomrule
\end{tabular}
\end{table}

Bei $n = 5$ ist das TI mehr als doppelt so breit wie bei
$n \to \infty$. Die Investition in mehr Messwerte lohnt sich
besonders im Bereich $n = 5$--$30$.


% ══════════════════════════════════════════════════════════════
\section{Ausblick: Zensierte Daten}
\label{sec:zensiert}
% ══════════════════════════════════════════════════════════════

In der Praxis der Lebensdauerprüfung werden Tests häufig
abgebrochen, bevor alle Prüflinge ausgefallen sind. Die
resultierenden Daten sind \emph{zensiert}.

\subsection{Zensierungsarten}

\begin{description}[nosep,leftmargin=*,font=\bfseries]
    \item[Typ-I-Zensierung (zeitlich):] Der Test endet nach
        einer festen Zeit $T$. Prüflinge, die bis $T$ nicht
        ausgefallen sind, liefern nur die Information
        ,,Lebensdauer $> T$''.
    \item[Typ-II-Zensierung (anzahlbasiert):] Der Test endet
        nach dem $r$-ten Ausfall. Die Lebensdauern der
        verbleibenden $n - r$ Prüflinge sind rechtszensiert.
    \item[Zufällige Zensierung:] Prüflinge fallen zu
        verschiedenen Zeitpunkten aus dem Test (z.\,B.\ durch
        Bruch, Verlust). Häufigster Fall in der Praxis.
\end{description}

\subsection{Auswirkung auf Toleranzintervalle}

Zensierte Daten erfordern Anpassungen aller Verfahren:
\begin{itemize}[nosep,leftmargin=*]
    \item \textbf{Verteilungsfrei}: Nicht direkt anwendbar,
          da $x_{(n)}$ nicht beobachtet wird
    \item \textbf{Normal/Lognormal}: $\xbar$ und $s$ müssen
          durch zensierungsbereinigte Schätzer ersetzt werden
          (z.\,B.\ Kaplan-Meier für den Median)
    \item \textbf{Weibull}: Die MLE-Likelihood muss um
          Zensierungsterme erweitert werden:
\end{itemize}

\begin{equation}
    \ell(c, \lambda) = \sum_{i \in \mathcal{F}} \ln f(t_i)
    + \sum_{i \in \mathcal{C}} \ln S(t_i)
    \label{eq:censored_lik}
\end{equation}

wobei $\mathcal{F}$ die Menge der beobachteten Ausfälle,
$\mathcal{C}$ die Menge der zensierten Beobachtungen und
$S(t) = 1 - F(t)$ die Überlebensfunktion ist.

Für Weibull ergibt sich konkret:
\begin{equation}
    S(t) = \exp\!\left[-\left(\frac{t}{\lambda}\right)^c\right]
\end{equation}

\hinweis{Die aktuelle Implementierung unterstützt noch keine
zensierten Daten. Für diesen Anwendungsfall empfiehlt sich die
Python-Bibliothek \texttt{lifelines} oder \texttt{reliability},
die zensierungsbereinigte MLE und Bootstrap implementieren.}


% ══════════════════════════════════════════════════════════════
\section{Fehlerquellen und häufige Missverständnisse}
\label{sec:fehler}
% ══════════════════════════════════════════════════════════════

\subsection{Verwechslung von Konfidenz und Abdeckung}

Ein 95/95-TI wird gelegentlich als ,,95\%-Intervall''
abgekürzt -- aber \emph{welche} 95\%? Es ist essenziell,
beide Parameter anzugeben. Ein 99/50-TI (99\% Abdeckung,
50\% Konfidenz) ist eine völlig andere Aussage als ein
50/99-TI (50\% Abdeckung, 99\% Konfidenz).

\subsection{Normalitätsannahme bei kleinem $n$}

Bei $n = 10$ hat der Shapiro-Wilk-Test eine Macht von
typischerweise 30--50\% gegen moderate Abweichungen von der
Normalverteilung. Das bedeutet: Selbst wenn die Daten
\emph{nicht} normalverteilt sind, wird $H_0$ in mehr als
der Hälfte der Fälle \emph{nicht} abgelehnt. Der Test
,,bestätigt'' also nicht die Normalität -- er kann sie nur
bei starken Abweichungen ablehnen.

Konsequenz: Bei kleinem $n$ sollte domänenspezifisches
Wissen über die Verteilung einfließen, nicht blind dem
Testergebnis vertraut werden.

\subsection{Multiple Testproblematik}

Der Entscheidungsalgorithmus testet sequenziell auf Normal-,
Lognormal- und Weibull-Verteilung. Jeder Test hat eine
Fehlerwahrscheinlichkeit $\alpha_s$. Die Gesamtwahrscheinlichkeit,
mindestens eine falsche Akzeptanz zu produzieren, ist
\emph{nicht} $\alpha_s$, sondern kann bis zu
$1 - (1-\alpha_s)^3 \approx 0.143$ für $\alpha_s = 0.05$ betragen.

In der Praxis ist dies kein gravierendes Problem, weil die
Tests sequenziell (nicht parallel) durchgeführt werden und
bei der ersten Akzeptanz abgebrochen wird. Die Prioritätsreihenfolge
(Normal vor Lognormal vor Weibull) stellt sicher, dass das
einfachste passende Modell gewählt wird.


% ══════════════════════════════════════════════════════════════
\section{Referenzen und weiterführende Literatur}
\label{sec:referenzen}
% ══════════════════════════════════════════════════════════════

\begin{enumerate}[nosep,leftmargin=*,label={[\arabic*]}]
    \item \textsc{Meeker}, W.\,Q., \textsc{Hahn}, G.\,J. und
          \textsc{Escobar}, L.\,A.:
          \emph{Statistical Intervals: A Guide for Practitioners
          and Researchers}. 2.~Aufl., Wiley, 2017.

    \item \textsc{Howe}, W.\,G.:
          Two-sided tolerance limits for normal populations --
          some improvements.
          \emph{J.\ Amer.\ Statist.\ Assoc.}, 64:610--620, 1969.

    \item ISO~16269-6:2014:
          \emph{Statistical interpretation of data -- Part~6:
          Determination of statistical tolerance intervals}.

    \item \textsc{Shapiro}, S.\,S. und \textsc{Wilk}, M.\,B.:
          An analysis of variance test for normality.
          \emph{Biometrika}, 52(3/4):591--611, 1965.

    \item \textsc{Guenther}, W.\,C.:
          Tolerance intervals for univariate distributions.
          \emph{Naval Res.\ Logist.\ Quart.}, 19(2):309--333, 1972.

    \item \textsc{Krishnamoorthy}, K. und \textsc{Mathew}, T.:
          \emph{Statistical Tolerance Regions: Theory, Applications,
          and Computation}. Wiley, 2009.

    \item \textsc{Lawless}, J.\,F.:
          \emph{Statistical Models and Methods for Lifetime Data}.
          2.~Aufl., Wiley, 2003.

    \item NIST/SEMATECH e-Handbook of Statistical Methods:
          \emph{Tolerance Intervals}.
          \url{https://www.itl.nist.gov/div898/handbook/prc/section2/prc263.htm}

    \item \textsc{Wilks}, S.\,S.:
          Determination of sample sizes for setting tolerance limits.
          \emph{Ann.\ Math.\ Statist.}, 12(1):91--96, 1941.
\end{enumerate}

\vfill

\begin{center}
\small\textit{Dokument erstellt im Februar 2026.\\
Implementierung in Python 3.12+ mit NumPy, SciPy und PySide6.}
\end{center}

\newpage

% ══════════════════════════════════════════════════════════════
\appendix
\section{Formelsammlung}
\label{app:formeln}
% ══════════════════════════════════════════════════════════════

Zum schnellen Nachschlagen -- alle zentralen Formeln auf einen
Blick.

\subsection{Verteilungsfreies TI}

\textbf{Allgemein} ($\ell, u$ beliebig):
\begin{align}
    &\text{Konfidenz:} \quad \mathrm{CPTI} = F_{\mathrm{Binom}}(d-1;\, n,\, p) \\
    &\text{wobei } d = u - \ell \text{ (Span)}
\end{align}

\textbf{Optimiert} (Satz~\ref{satz:opt_df}):
\begin{align}
    &d^* = \min\{d : F_{\mathrm{Binom}}(d-1;\, n,\, p) \geq 1-\alpha\} \\
    &\ell^* = \arg\min_\ell \big(x_{(\ell+d^*)} - x_{(\ell)}\big) \\
    &\text{Intervall:} \quad [x_{(\ell^*)},\; x_{(\ell^*+d^*)}]
\end{align}

\textbf{Spezialfall} ($\ell=1, u=n$):
\begin{align}
    &\text{Konfidenz:} \quad C = 1 - np^{n-1} + (n-1)p^n \\
    &\text{Min.\ } n: \quad 1 - np^{n-1} + (n-1)p^n \geq 1 - \alpha
\end{align}

\textbf{Einseitig:}
\begin{align}
    &\text{Konfidenz:} \quad C = 1 - p^n \\
    &\text{Min.\ } n: \quad n \geq \frac{\ln\alpha}{\ln p}
\end{align}

\subsection{Normal-TI}

\textbf{Zweiseitig:}
\begin{align}
    &\text{Intervall:} \quad [\xbar - ks,\; \xbar + ks] \\
    &k = z_{(1+p)/2} \cdot
    \sqrt{\frac{(n-1)(1+1/n)}{\chi^2_{1-\alpha,\,n-1}}}
\end{align}

\textbf{Einseitig:}
\begin{align}
    &\text{Intervall:} \quad (-\infty,\; \xbar + ks] \\
    &k = \frac{t_{n-1,\,1-\alpha,\,z_p\sqrt{n}}}{\sqrt{n}}
\end{align}

\subsection{Lognormal-TI}

Setze $y_i = \ln x_i$, berechne $\bar{y}$ und $s_y$:
\begin{equation}
    \Big[\exp(\bar{y} - k \cdot s_y),\;\;
          \exp(\bar{y} + k \cdot s_y)\Big]
\end{equation}
mit $k$ wie im Normalfall.

\subsection{Weibull-TI}

Weibull-Dichte und Quantilfunktion:
\begin{align}
    f(x) &= \frac{c}{\lambda}
    \left(\frac{x}{\lambda}\right)^{\!c-1}
    \exp\!\left[-\left(\frac{x}{\lambda}\right)^{\!c}\right] \\[4pt]
    F^{-1}(q) &= \lambda\,\big(-\ln(1-q)\big)^{1/c}
\end{align}

Bootstrap-Konfidenzgrenzen ($B$ Replikationen):
\begin{align}
    L &= Q_{\alpha/2}\!\left(\big\{
        F^{-1}\!\big(\tfrac{1-p}{2};\, \hat{c}^{(b)}, \hat{\lambda}^{(b)}\big)
    \big\}_{b=1}^B\right) \\
    U &= Q_{1-\alpha/2}\!\left(\big\{
        F^{-1}\!\big(\tfrac{1+p}{2};\, \hat{c}^{(b)}, \hat{\lambda}^{(b)}\big)
    \big\}_{b=1}^B\right)
\end{align}

\subsection{Shapiro-Wilk-Teststatistik}

\begin{equation}
    W = \frac{\left(\sum_{i=1}^n a_i\, x_{(i)}\right)^2}
    {\sum_{i=1}^n (x_i - \xbar)^2}
    \quad\text{mit}\quad
    \mathbf{a} = \frac{\mathbf{V}^{-1}\mathbf{m}}
    {\|\mathbf{V}^{-1}\mathbf{m}\|}
\end{equation}

\subsection{KS-Teststatistik (Weibull-GoF)}

\begin{equation}
    D_n = \sup_x \big|F_n(x) - F(x;\hat{c},\hat{\lambda})\big|
\end{equation}

Monte-Carlo-$p$-Wert ($N_{\text{MC}}$ Replikationen):
\begin{equation}
    p_{\text{MC}} = \frac{1}{N_{\text{MC}}}
    \sum_{j=1}^{N_{\text{MC}}} \mathbf{1}\big[D_n^{(j)} \geq D_n\big]
\end{equation}

\subsection{Prozessfähigkeit}

\begin{align}
    C_p &= \frac{\text{USL} - \text{LSL}}{6\sigma} \\[4pt]
    C_{pk} &= \min\!\left(
        \frac{\text{USL} - \mu}{3\sigma},\;
        \frac{\mu - \text{LSL}}{3\sigma}
    \right)
\end{align}

TI-basierte Spezifikationsprüfung:
\begin{equation}
    \text{LSL} \leq \xbar - ks
    \quad\text{und}\quad
    \xbar + ks \leq \text{USL}
\end{equation}

% ══════════════════════════════════════════════════════════════
\end{document}
% ══════════════════════════════════════════════════════════════
